
%%%%%%%%%%%%%%%%%%%%%%%%%%%%%%%%%%%%%%%%%%%%%%%%%%%%%%%%%%%%%%%%%%%%%%%%%%%%%%%%%%%%%%%%%%%%%%%%%%% 
% KAPITOLA 5
%%%%%%%%%%%%%%%%%%%%%%%%%%%%%%%%%%%%%%%%%%%%%%%%%%%%%%%%%%%%%%%%%%%%%%%%%%%%%%%%%%%%%%%%%%%%%%%%%%
\chapter{Návrh}
\label{ch:kap5}
Táto kapitola sa zaoberá vysokoúrovňovým návrhom nástroja na automatizované generovanie receptov pre nástroj Dataster. Následne približuje celkovú architektúru (Sekcia~\ref{sub:architektura}), komponenty návrhu a tok dát medzi nimi (Sekcia~\ref{sub:fazy_generovania}). LLM tvorí jadro návrhového riešenia a preto je pozornosť venovaná spôsobu využitia veľkých jazykových modelov (Sekcia~\ref{sub:vyuzitie_vlk_jaz_modelov}). Pokračuje princípom validácie a opravy výstupov (Sekcia~\ref{sub:validacia_a_kontrola}) a na záver opisuje rozdelenie práce (Sekcia~\ref{sub:rozdelenie_prace}). Konkrétna implementácia návrhu je obsiahnutá v Kapitole~\ref{ch:kap6}.

\section{Architektúra}

\label{sub:architektura}

Vstupom nástroja sú identifikačné údaje pre pripojenie k relačnej databáze. Na začiatku sa program pripojí k databáze, extrahuje tabuľky, informácie o nich, teda stĺpce, ich dátové typy, primárne a cudzie kľúče. Po získaní týchto informácií sa z tabuľky vytiahnu rôzne vzorky krajných hodnôt z oboch koncov usporiadania, ktoré vymedzujú rozsah hodnôt stĺpca, čo jazykovému modelu umožní lepšie pochopiť charakter jednotlivých stĺpcov.



V ďalšej fáze LLM ako prvé vygeneruje sekciu \texttt{types}, keďže na tejto sekcii závisia ostatné sekcie YAML. Následne sa vygenerujú ďalšie sekcie. Sekcie schema, generate a constraints sa generujú nezávisle, každá samostatným volaním modelu s types sekciou ako vstupom. 



Surové výstupy jednotlivých fáz sa následne poskladajú do jedného predbežného YAML, ktorý je validovaný a v prípade chýb opravený. Zložený recept prechádza automatickými opravami, následne jednou validáciou a uloží sa.

\begin{figure}[h]
    \centering
    \includegraphics[width=1\linewidth]{obrazky-figures/diagram1.pdf}
    \caption{Architektúra nástroja}
    \label{fig:diagram1}
\end{figure}

\section{Fázy generovania}

\label{sub:fazy_generovania}

Samotné generovanie sekcií je rozdelené do niekoľkých fáz, hlavne z dôvodu izolácie chýb. Pri zlyhaní nie sú ostatné sekcie stratené a chyba je priraditeľná fáze, v ktorej vznikla. Malé modely majú problém s dodržaním konzistentného formátu výstupu, pričom použitie few-shot príkladov tento problém pomáha zmierniť~\cite{Pysmennyi2025}. Rozdelenie úloh šetrí pamäť. Každý krok spracováva len nutné informácie, čo je pri 8 GB VRAM zásadná výhoda.


Pre \textbf{fázu types} vstup obsahuje informácie o primárnom kľúči, dátových typoch stĺpcov, ako aj niekoľko vzoriek hodnôt. Ide o interpretačne najnáročnejšiu fázu a jej kvalita vygenerovaného výstupu priamo ovplyvňuje ostatné fázy, preto sa jej návrhu venuje najväčšia pozornosť. Databázové vzorky vidí iba fáza types, ostatné fázy pracujú výlučne s jej výstupom.


\textbf{Fáza schema} prevezme vygenerovanú sekciu types. Jej úlohou je mapovanie stĺpcov a vzťahov cudzích kľúčov. Všetky cudzie kľúče vytiahne program priamo z databázy a do schémy ich vloží automaticky. Model sa používa najmä na návrh prepojení a na odhaľovanie vzťahov medzi tabuľkami, ktoré nie sú v databáze explicitne zapísané a treba ich odhadnúť podľa názvov a hodnôt.


Pri \textbf{fáze generate} sa určujú jednotlivé počty generovaných riadkov pre každú tabuľku. Logický počet sa odvíja od rozsahu hodnôt primárneho kľúča.


Ako posledná \textbf{fáza} je \textbf{constraints}, ktorá podľa vstupu types odvodzuje logické obmedzenia hodnôt, ktoré nie sú zachytené v schéme databázy, ale vyplývajú z kontextu hodnôt, ako napríklad rozsahy veku alebo nezáporné ceny. Tieto obmedzenia nie je možné mechanicky vyčítať, model ich musí vyvodiť zo sekcie types.


Toto poradie je dôležité, pretože types je spoločným vstupom ostatných troch generovaných sekcií. Types definuje množinu tabuliek, stĺpcov a ich typov, na ktorú sa ostatné sekcie odkazujú. Ostatné fázy sú navzájom nezávislé a preto sa volajú oddelene, chyba v jednej z týchto častí neruší ostatné. 

\section{Využitie veľkých jazykových modelov (LLM)}

\label{sub:vyuzitie_vlk_jaz_modelov}

Jazykový model v návrhu nefunguje ako autonómny agent, ktorý sám rozhoduje o ďalšom kroku, poradie fáz vrátane validácie a opráv sú pevne riadené programom. AI agenti sú síce populárni v oblasti zabezpečovania kvality softvéru~\cite{Pysmennyi2025}, ale pre tento program to bolo vyhodnotené ako nevhodné hneď z niekoľkých dôvodov. Po prvé, experimentálne vyhodnotenie vyžaduje reprodukovateľnosť, rovnaký vstup musí prejsť tou istou cestou. S agentom pribudne ďalšia neistota, pretože pri každom spustení sa môže líšiť nielen odpoveď modelu, ale aj počet a poradie krokov. Po druhé, spoľahlivé plánovanie je zdokumentovaným problémom aj pri väčších modeloch~\cite{Valmeekam2023}, pri modeloch s okolo 7 miliardami parametrov, ktoré si vynucuje dostupná pamäť (8 GB VRAM), ho nemožno preto predpokladať. Po tretie, keď má každá fáza svoje pevné miesto, zlyhanie je priraditeľné k fáze, v ktorej sa objavilo, pri slučke tvorenej agentom sa chyba stráca v trajektórii.


Modely tak rozhodujú výhradne o obsahu generovaných sekcií, nie o štruktúre procesu. Všetko, čo je možné vykonať algoritmicky, vykonáva kód.

\section{Validácia a oprava}

\label{sub:validacia_a_kontrola}

Priebeh validácie zloženého receptu znázorňuje obrázok~\ref{fig:diagram2}, teda pri validácii môžu nastať dva prípady. Buď je vygenerovaný YAML bez chýb a uloží sa, alebo prechádza kontrolou a radom opráv, pričom sa odchytené chyby zaznamenajú a recept sa uloží. Tento prístup bol zvolený zámerne, aby použivateľ mohol recept opraviť alebo nástroj spustiť znova. Zároveň to umožňuje vyhodnocovať aj recepty, ktoré validáciou neprešli (Kapitola~\ref{ch:kap7}). V prípade, že kód narazí na vážnu chybu, kód končí s chybovým hlásením.  

\begin{figure}[h]
    \centering
    \includegraphics[width=0.75\linewidth]{obrazky-figures/diagram2.pdf}
    \caption{Schéma automatickej opravy a validácie zloženého receptu}
    \label{fig:diagram2}
\end{figure}


Výstupy jazykových modelov sú formátovo nestabilné a menšie modely majú zdokumentované problémy s dodržaním konzistentného výstupného formátu~\cite{Pysmennyi2025}. Preto bola automatická oprava a kontrola výstupov navrhnutá ako integrálna súčasť procesu, nie ako dodatočná kontrola.

\section{Rozdelenie práce}

\label{sub:rozdelenie_prace}

Využitie LLM pre každú časť procesu, od extrakcie schémy cez generovanie až po opravy by bolo nielen zbytočné, ale aj pomalé. Návrh túto prácu delí medzi LLM a samostatný kód podľa povahy úlohy. 


LLM sa stará o generovanie YAML častí ako \texttt{types}, \texttt{schema}, \texttt{generate} a \texttt{constraints}, ale z generovania vynecháva sekciu \texttt{include} a SQL požiadavky. LLM dostáva úlohy, pri ktorých je dôležitá interpretácia a neexistuje presný algoritmus.


Samotný kód dostáva úlohy na mechanickej báze, kde je nutný presný postup a odchýlka od neho je chybou. Stará sa o pripojenie k databáze, extrakciu schémy a vzoriek dát ako aj poskladanie týchto informácií do jedného balíka, ktorý je poslaný do LLM. Ako posledné si prevezme vygenerovaný výsledok, ten skontroluje, opraví a uloží. 


V počiatočných fázach návrhu bol testovaný aj variant fázy, ktorý sám opravoval chybný recept, lenže opravy boli pomalé a nespoľahlivé a nakoniec sa oprava kódom ukázala ako lepšia voľba. 


Pôvodný návrh zo zadania počítal s LLM aj pri generovaní SQL dotazov ako aj generovaní YAML include sekcie. Obe úlohy sa ale ukázali byť plne programovateľné a keďže LLM nepridávalo hodnotu ani jednej zo spomenutých častí a vytváral zbytočné chyby, implementácia týchto častí bola zverená kódu.


